\section{Introduction}
\label{sec:introduction}

% Para 1 -- The converging attack surface
In 2024, researchers discovered six zero-day vulnerabilities in each of 16 EV charging management systems~\cite{sarieddine2024ocpp}, one symptom of a broader shift toward software-managed power grids where cloud APIs control charging fleets, DER setpoints, and demand response across thousands of devices.
Adversaries who compromise an EV or DER management platform gain programmatic control over aggregate load, enabling coordinated load-altering attacks that can propagate from distribution feeders into the transmission system.
In the DER domain, the Forescout SUN:DOWN report identified 129 CVEs across major solar inverter management platforms exposing over 1~TW of capacity to remote control, and Kuroptev et al.~\cite{kuroptev2026evgrid} demonstrated using real German charging station data that compromising just two Charging Station Operators suffices to overload transmission grid branches beyond N-1 security margins.
As these programmatic attack surfaces proliferate, the question of how to evaluate grid defenses against adversaries who exploit them becomes increasingly urgent.

% Para 2 -- The evaluation problem and EVG concept
The smart grid security community has developed a rich ecosystem of co-simulation frameworks and intrusion detection testbeds~\cite{cintuglu2017testbeds, smadi2021testbeds}, but most evaluation methodologies still rely on static datasets---widely used traces such as SWaT~\cite{goh2017swat}, WADI~\cite{ahmed2017wadi}, EPIC~\cite{adepu2019epic}, and HAI~\cite{shin2020hai} capture fixed attack trajectories under specific operating conditions---or on a small number of pre-scripted attack scenarios as in GridAttackSim~\cite{le2020gridattacksim}.
Static traces and scripts cannot capture the action--reaction dynamics between adaptive attackers and responsive defenses. When a controller sheds load in response to an attack, a static attacker continues its pre-programmed sequence regardless, while an adaptive attacker can observe the defensive response and adjust its strategy in real time.
We term the resulting discrepancy the \emph{evaluation validity gap} (EVG): the ratio of a defense's stress under an adaptive attacker to its stress under a comparable static attacker, where EVG $> 1.0$ indicates that static evaluations underestimate adversarial risk.
This phenomenon has been extensively documented in adversarial machine learning (Carlini et al.~\cite{carlini2019evaluating} established best practices for adaptive evaluation, and Nasr et al.~\cite{nasr2025attacker} recently demonstrated that 12 LLM defenses reporting near-zero attack success rates under static evaluation were bypassed at over 90\% success under adaptive attacks) but has not been formally studied or quantified in smart grid security evaluation.
This paper asks: does a measurable evaluation validity gap exist in smart grid security, and if so, what attacker capabilities drive it?

% Para 3 -- Our approach: LLMs as adaptive attackers
To answer this question, we introduce LLM-GridEval, a framework that embeds tool-using LLM agents as adaptive attackers inside a HELICS-based transmission--distribution co-simulation~\cite{hardy2024helics}, connecting them to the grid through a schema-constrained catalog of attack primitives that ensures all actions are physically realizable.
We use tool-using LLMs as a practical way to instantiate adaptive attackers without environment-specific training, unlike reinforcement learning approaches that demand scenario-specific reward shaping and state representations~\cite{sahani2023mlids}.
A critical architectural choice is the schema-constrained primitive interface, which sidesteps the LLM code-generation reliability problem identified by Sen et al.~\cite{sen2024aiattacker}, who found that current LLMs produce only ``partially usable'' attack code for grid systems; by constraining the LLM to \emph{select} validated actions rather than \emph{generate} arbitrary commands, we make LLM-driven grid attacks feasible and reproducible.
We emphasize that LLMs are not proposed as optimal attackers but as a practical, explainable, and generalizable approach to adaptive adversary modeling that complements existing static and RL-based methods.

% Para 4 -- Contributions
We make the following contributions:
\begin{enumerate}[leftmargin=*]
\item \textbf{Framework.} We present LLM-GridEval, to our knowledge the first framework coupling tool-using LLM agents with physics-accurate T\&D co-simulation via schema-constrained attack primitives for adversarial security evaluation. The framework's layered architecture enables controlled experimentation by allowing researchers to swap attacker implementations (scripted, RL-based, or LLM-based) while holding the co-simulation and tool interface fixed.

\item \textbf{Empirical results.} We conduct a 37-experiment campaign across three grid operating points (low, medium, and high base load) comparing three attacker variants---a random static baseline, an LLM attacker with timing intelligence only, and an LLM attacker with the full AI-V2 prompt package (timing, domain knowledge, diversification guidance, attack history, and fixed max-power behavior)---against a progressive-shedding controller. The strategy-aware attacker achieves EVG $= 2.25$--$2.40\times$ at operating points where the controller can respond ($p < 0.05$, $d = 1.85$--$4.95$), while EVG equals $1.0\times$ at a ceiling operating point where the controller is overwhelmed.

\item \textbf{Insight on attacker capabilities.} We show that the AI-V2 prompt package, combining domain knowledge, diversification guidance, and dynamic attack history, substantially outperforms timing-only intelligence, while timing intelligence alone appears counterproductive at low-load conditions in this setup (EVG $= 0.80\times$), partly due to safety refusals and single-target fixation, providing guidance for both prompt engineering of CPS red-teaming agents and the design of adaptive-aware defense evaluations.
\end{enumerate}
